\documentclass[12pt,a4paper]{article}

\usepackage[margin=2.5cm]{geometry}
\usepackage{graphicx}
\usepackage{booktabs}
\usepackage{array}
\usepackage{amsmath}
\usepackage{siunitx}
\usepackage{natbib}
\usepackage{caption}
\usepackage{subcaption}
\usepackage{setspace}
\usepackage{lineno}
\usepackage{url}
\usepackage[hidelinks]{hyperref}

\graphicspath{
  {figures/}
  {D:/Data_Analysis_R/Gas_Concentration_Emission_R/workflows/mapping_high_resolution/plots/manuscript1/}
}

\sisetup{
  separate-uncertainty = true,
  per-mode = symbol,
  range-phrase = {--},
  range-units = single
}

\onehalfspacing
\linenumbers

\title{Spatial heterogeneity of carbon dioxide, methane, and ammonia in a naturally ventilated dairy barn: high-resolution mapping and implications for sampling design}

\author{
Harsh Sahu$^{a,b}$,
Long Chen$^{a}$,
Thomas Amon$^{a,b}$,
David Janke$^{a,*}$
}

\date{}

\begin{document}

\maketitle

\begin{center}
\small
$^{a}$Leibniz Institute for Agricultural Engineering and Bioeconomy (ATB),
Max-Eyth-Allee 100, 14469 Potsdam, Germany\\
$^{b}$Institute of Animal Hygiene and Environmental Health, Department of
Veterinary Medicine, Freie Universit\"at Berlin, Germany\\
$^{*}$Corresponding author: \href{mailto:djanke@atb-potsdam.de}{djanke@atb-potsdam.de}
\end{center}

\begin{abstract}
Homogeneous gas mixing is commonly assumed when concentrations are used to
characterise air quality or calculate emissions from naturally ventilated
dairy barns, but the magnitude and spatial structure of departures from this
assumption remain poorly resolved. Carbon dioxide (\(\mathrm{CO_2}\)),
methane (\(\mathrm{CH_4}\)), and ammonia (\(\mathrm{NH_3}\)) concentrations
were therefore mapped at 51 internal locations arranged as 17 vertical
triplets. Two cavity ring-down spectrometers and two Fourier-transform infrared
spectrometers were operated at a temporal resolution of four minutes; the
Fourier-transform infrared measurements were harmonised to the cavity
ring-down spectroscopy scale using a preliminary 24 h co-located comparison.
All gases exhibited structured vertical and horizontal heterogeneity. Mean
concentrations decreased from top to middle to bottom, but the magnitude of
the vertical gradient differed among gases and horizontal positions. The
\(\mathrm{CH_4}/\mathrm{CO_2}\) and
\(\mathrm{NH_3}/\mathrm{CO_2}\) ratios also decreased with height group from
top to bottom, showing that the observed structure was not explained solely by
uniform dilution or concentration. As a practical consequence of this spatial
analysis, alternative height configurations were evaluated in 743 complete
two-hour blocks. Top-only and bottom-only measurements produced systematic,
gas-specific biases, whereas a spatially distributed top-and-bottom network
reproduced the full-density mean with Lin's concordance correlation
coefficients of 0.999 and mean relative biases below 0.4\%. The primary finding
is that gas concentrations within the barn were spatially structured rather
than homogeneous; understanding this structure subsequently enabled the
network to be reduced from 51 to 34 locations without materially changing the
barn mean.
\end{abstract}

\noindent\textbf{Keywords:}
air quality; dairy housing; horizontal gradient; mixing ratio; spatial
variability; vertical profile

\section*{Nomenclature}

\begin{tabular}{>{\raggedright\arraybackslash}p{2.2cm}
                >{\raggedright\arraybackslash}p{10.8cm}}
\toprule
Symbol & Definition\\
\midrule
\(\mathrm{CH_4}\) & methane\\
\(\mathrm{CO_2}\) & carbon dioxide\\
CCC & Lin's concordance correlation coefficient\\
CRDS & cavity ring-down spectroscopy\\
FTIR & Fourier-transform infrared spectroscopy\\
\(\mathrm{H_2O}\) & water vapour\\
\(\mathrm{NH_3}\) & ammonia\\
NVDB & naturally ventilated dairy barn\\
RE & relative deviation from the full-density estimate, \%\\
RMSE & root mean square error\\
\(s\) & south-background sampling location\\
SE & standard error\\
SP & sampling point\\
\bottomrule
\end{tabular}

\section{Introduction}

Concentrations measured in naturally ventilated livestock buildings are
affected by unsteady wind-driven exchange, local sources, animal activity,
building geometry, and imperfect internal mixing. Consequently, a
concentration recorded at one position cannot automatically be assumed to
represent the complete barn airspace. This representativeness is important
both for characterising animal exposure and for calculating gaseous emissions
with indirect balance methods.

Spatial variability has been demonstrated in both the horizontal and vertical
directions. In a naturally ventilated dairy cow barn, \citet{Mendes2015}
reported height-dependent concentrations of \(\mathrm{NH_3}\),
\(\mathrm{CO_2}\), and a released tracer gas, with mixing ratios becoming more
stable above the animal-occupied zone. Measurements at three vertical
locations in an open dairy barn also showed gas-specific profiles and only
moderate correlations between some heights \citep{Durso2024}. Horizontal
representativeness is similarly dependent on both the number and placement of
sampling points. \citet{Durso2022} found that variability decreased when the
number of horizontally distributed locations increased and that the distance
between locations affected the resulting mean.

Increasing sampling density can reduce uncertainty, but it also increases
equipment requirements, tubing, switching-cycle duration, data-processing
effort, and the probability of sampling-system failure. In a naturally
ventilated animal mock-up building, \citet{DeVogeleer2017} showed that the
spatial pattern of sensors could be as influential as the number of sensors.
Systematic reduction of a high-resolution wind-tunnel measurement grid
similarly showed that concentration and velocity sensors positioned at only
one height could produce large emission errors, whereas vertically composite
sampling improved representativeness \citep{Janke2020}. More recently,
information-based sensor selection has been shown to retain most predictive
performance with substantially fewer strategically selected locations
\citep{Chen2026}.

Available studies therefore indicate that dense sampling is valuable, but do
not establish how a full-scale, multi-height concentration network can be
reduced without losing the barn mean and its spatial heterogeneity under field
conditions. The objective of this study was to quantify the spatial
heterogeneity of \(\mathrm{CO_2}\), \(\mathrm{CH_4}\), and
\(\mathrm{NH_3}\) in a naturally ventilated dairy barn and to test the
hypothesis that a spatially distributed top-and-bottom configuration would
reproduce the concentration statistics obtained from a 51-location,
three-height reference design.

\section{Materials and methods}

\subsection{Experimental barn}

Measurements were conducted in the experimental dairy barn of the
Lehr- und Versuchsanstalt f\"ur Tierzucht und Tierhaltung, Gross Kreutz,
Brandenburg, Germany. The naturally ventilated barn was \SI{38.88}{m} long and
\SI{17.65}{m} wide. The asymmetrical fibre-cement roof reached approximately
\SI{6.2}{m} at the ridge and approximately \SI{3.6}{m} at the sidewalls. The
internal airspace was approximately \SI{4529}{m^3}. The loose-housing barn was
designed for 58 dairy cows and contained littered cubicles, concrete
walkways, a manure scraper, and adjacent milking and young-stock areas.

\subsection{Sampling layout}

Fifty-one internal SPs were distributed over 17 horizontal positions
(Fig.~\ref{fig:layout}). Three vertically aligned locations were installed at
each position: a top location \SI{60}{cm} below the roof, a middle location
\SI{360}{cm} above the floor, and a bottom location \SI{260}{cm} above the
floor. The location numbers were assigned consecutively within each triplet.
Thus, locations \(1,4,\ldots,49\) represented the top group; locations
\(2,5,\ldots,50\) represented the middle group; and locations
\(3,6,\ldots,51\) represented the bottom group. A single south-background SP,
denoted \(s\), was treated conceptually as location 52 but was excluded from
all analyses of internal spatial heterogeneity.

\begin{figure}[htbp]
\centering
\begin{subfigure}{\textwidth}
  \centering
  \includegraphics[width=\textwidth]{Fig1a_campaign1_sampling_layout_no_ringline.png}
  \caption{Plan view.}
\end{subfigure}

\vspace{0.5em}
\begin{subfigure}{0.92\textwidth}
  \centering
  \includegraphics[width=\textwidth]{Fig1b_barn_cross_section.png}
  \caption{Cross-sectional representation of the sampling heights and fans.}
\end{subfigure}
\caption{Experimental barn and high-resolution sampling design. Orange symbols
indicate locations \SI{60}{cm} below the roof, green symbols indicate
locations \SI{360}{cm} above the floor, and blue symbols indicate locations
\SI{260}{cm} above the floor. The numbered internal locations comprise 17
vertical triplets. The yellow point labelled Background represents the
south-background location \(s\), which was retained in the dataset but
excluded from internal spatial comparisons.}
\label{fig:layout}
\end{figure}

\subsection{Gas analysers and preliminary field comparison}

Concentrations of \(\mathrm{CO_2}\), \(\mathrm{CH_4}\),
\(\mathrm{NH_3}\), and water vapour were measured using two CRDS analysers
(CRDS8 and CRDS9) and two FTIR analysers (FTIR1 and FTIR2). Before the first
mapping campaign, all four analysers were operated for 24 h at a co-located
central barn position. Two polytetrafluoroethylene sampling tubes were used:
one was divided with a T-piece between the two CRDS analysers and the other
was divided between the two FTIR analysers. The FTIR instruments measured
\(\mathrm{CO_2}\) and \(\mathrm{CH_4}\) concentrations that were 6\% higher
than those measured by CRDS, while FTIR \(\mathrm{NH_3}\) was 9\% higher.
Accordingly, FTIR measurements from Campaign 1 were harmonised to the CRDS
scale as

\begin{align}
C_{\mathrm{CO_2,corr}} &= C_{\mathrm{CO_2,raw}}/1.06,\\
C_{\mathrm{CH_4,corr}} &= C_{\mathrm{CH_4,raw}}/1.06,\\
C_{\mathrm{NH_3,corr}} &= C_{\mathrm{NH_3,raw}}/1.09.
\end{align}

CRDS measurements and all water-vapour measurements were retained without
correction. Both raw and corrected concentrations were preserved in the
analytical dataset. Before and after the field comparison, all analysers were
tested with certified cylinders containing approximately \SI{700}{ppm}
\(\mathrm{CO_2}\), \SI{7}{ppm} \(\mathrm{CH_4}\), and \SI{5}{ppm}
\(\mathrm{NH_3}\); deviations remained within the manufacturer-stated
\(\pm 5\%\) accuracy. The field comparison was therefore used for
cross-platform harmonisation under representative sampling conditions.

\subsection{Measurement campaigns and data processing}

Campaign 1 data were recovered for June--August 2024. Simultaneous ordered
cycling by all four analysers was identified from 3 June to 16 July and again
from 1 to 26 August, with shorter interruptions retained as gaps. CRDS9
sampled locations 1--16, FTIR2 sampled locations 17--26,
CRDS8 sampled locations 27--42, and FTIR1 sampled locations 43--51 and \(s\).
The resulting network comprised 51 internal locations and the
south-background location. Gas measurements were recorded at a four-minute
temporal resolution. Each settled CRDS valve step was identified from the
ordered 1--16 sequence. The transition signal was excluded, the first
\SI{60}{s} was discarded, and the subsequent \SI{180}{s} was averaged.
Across June--August, CRDS9 contained 22,384 valid four-minute steps and two
isolated three-minute steps; CRDS8 contained 26,099 four-minute steps and one
three-minute step. For every three-minute step, \SI{120}{s} remained after
the common \SI{60}{s} flushing period. Single-position and irregular
sequences were retained in audit tables but excluded from the campaign data.

Campaign 2 was conducted from 16 November 2024 at 00:06:47 to 31 December
2024 at 23:57:10 using CRDS8 and CRDS9. The 17 middle locations were omitted
following the Campaign 1 density assessment. Locations 19 and 40 were also
excluded because of fan proximity and a sampling-line leakage problem,
respectively, resulting in 32 internal locations. Each CRDS sampling step
lasted \SI{240}{s}; the first \SI{60}{s} was discarded for flushing and the
remaining \SI{180}{s} was averaged.

Non-positive gas records were treated as missing in the manuscript analytical
dataset. Carbon dioxide values below \SI{300}{ppm} were also treated as
invalid transition or analyser records because they were physically
implausible for barn or ambient air. No row was deleted from the immutable
source dataset.

\begin{table}[htbp]
\centering
\caption{Measurement design and analytical records retained for the two
campaigns. The south-background location \(s\) was retained in the data but
excluded from the count of internal locations.}
\label{tab:campaigns}
\begingroup
\small
\setlength{\tabcolsep}{1.5pt}
\begin{tabular}{@{}c l
  >{\raggedright\arraybackslash}p{4.0cm}
  >{\raggedright\arraybackslash}p{2.7cm}
  r r@{}}
\toprule
Campaign & Period & Analysers & Vertical design & Internal SPs & Rows\\
\midrule
1 & Jun.--Aug. 2024 & CRDS8, CRDS9, FTIR1, FTIR2 &
top, middle, bottom & 51 & 111\,210\\
2 & 16 Nov.--31 Dec. 2024 & CRDS8, CRDS9 &
top and bottom & 32 & 20\,632\\
\bottomrule
\end{tabular}
\endgroup
\end{table}

\subsection{Gas ratios}

The relative mixing ratios of \(\mathrm{CH_4}\) and \(\mathrm{NH_3}\) to
\(\mathrm{CO_2}\) were calculated for positive, valid concentrations:

\begin{align}
R_{\mathrm{CH_4:CO_2}} &=
100\,C_{\mathrm{CH_4}}/C_{\mathrm{CO_2}},\\
R_{\mathrm{NH_3:CO_2}} &=
100\,C_{\mathrm{NH_3}}/C_{\mathrm{CO_2}},\\
R_{\mathrm{NH_3:CH_4}} &=
100\,C_{\mathrm{NH_3}}/C_{\mathrm{CH_4}}.
\end{align}

The ratios are dimensionless and are reported as percentages. Where
inferential modelling of ratios is subsequently performed, their logarithmic
forms, \(\log(C_{\mathrm{CH_4}}/C_{\mathrm{CO_2}})\) and
\(\log(C_{\mathrm{NH_3}}/C_{\mathrm{CO_2}})\), are preferred to avoid the
asymmetry of untransformed ratios.

\subsection{Evaluation of sampling-density reduction}

Daily means at each location were first calculated to provide independent
blocks for descriptive means and 95\% confidence intervals. For direct
comparison of sampling configurations, measurements were averaged into
two-hour blocks. A block was retained only when all 51 internal locations
were represented, yielding 743 complete blocks in the recovered dataset.

Seven vertical configurations were reconstructed from the Campaign 1 data:
top only, middle only, bottom only, top plus middle, middle plus bottom, top
plus bottom, and the full three-height design. The full 51-location mean was
used as the reference estimate within each block. For a reduced configuration
\(k\), the relative deviation was calculated as

\begin{equation}
RE_{k,t}=100\,
\frac{\overline{C}_{k,t}-\overline{C}_{51,t}}
     {\overline{C}_{51,t}}.
\end{equation}

The top-and-bottom configuration was selected a priori as the operational
reduced design because it retained both ends of the observed vertical
gradient. Its agreement with the full design was quantified using mean bias,
relative bias, RMSE, coefficient of determination, and Lin's CCC. Spatial
heterogeneity within each block was expressed as the standard deviation of
log-transformed location concentrations.

Following \citet{Chen2026}, the temporal informativeness of every location was
described using Shannon entropy. For each gas and ratio, location means from
complete two-hour blocks were divided into four equal-width bins. The
empirical probabilities \(p_b\) were used to calculate
\begin{equation}
H^*=-\frac{\sum_{b=1}^{4}p_b\log_2(p_b)}{\log_2(4)}.
\end{equation}
This was an empirical adaptation of the cross-case CFD entropy method: the
present cases were observed two-hour blocks rather than simulated ventilation
scenarios. Sensitivity to temporal aggregation was assessed at 2, 4, 8, and
24 h resolutions.

All cleaning, calculations, and figures were produced in R version 4.5.3
(R Foundation for Statistical Computing, Vienna, Austria). Data manipulation
was conducted using \texttt{data.table}, and figures were produced using
\texttt{ggplot2}. Top, middle, and bottom locations were displayed in orange,
green, and light blue, respectively.

\section{Results}

\subsection{High-resolution concentration patterns}

Marked spatial and vertical variation was detected during Campaign 1
(Fig.~\ref{fig:locations}). Across internal locations, the mean corrected
\(\mathrm{CO_2}\) concentrations were approximately \SI{616}{ppm} at the top,
\SI{611}{ppm} at the middle, and \SI{609}{ppm} at the bottom.
Corresponding \(\mathrm{CH_4}\) means were 18.84, 18.51, and
\SI{18.38}{ppm}, while \(\mathrm{NH_3}\) means were 2.05, 2.04, and
\SI{2.02}{ppm}. The direction and magnitude of the vertical gradient were not
uniform among horizontal positions. Particularly distinct values occurred at
locations 19 and 40, supporting their exclusion from the subsequent
operational design.

\begin{figure}[htbp]
\centering
\includegraphics[width=\textwidth]{Fig2_june_august_spatial_means_with_CV_v03.png}
\caption{Campaign 1 corrected concentration and gas-ratio means across the 17
horizontal positions. Symbol size represents the location-specific temporal
coefficient of variation, capped at 100\% for display. Orange, green, and
light-blue symbols denote top, middle, and bottom locations, respectively.
The south-background location \(s\) was excluded.}
\label{fig:locations}
\end{figure}

\subsection{Effect of vertical sampling configuration}

Use of one height alone produced systematic deviations from the full-density
mean (Fig.~\ref{fig:configurations}; Table~\ref{tab:configurations}).
Top-only sampling overestimated \(\mathrm{CO_2}\), \(\mathrm{CH_4}\), and
\(\mathrm{NH_3}\) by mean values of 0.68\%, 1.01\%, and 1.18\%,
respectively. Bottom-only sampling underestimated the same gases by 0.50\%,
0.78\%, and 1.40\%. The middle-only configuration was closer to the full mean,
but removed the vertical extremes required for evaluating heterogeneity.

Combining top and bottom locations substantially reduced configuration bias.
The mean absolute relative errors were 0.33\% for \(\mathrm{CO_2}\), 0.95\%
for \(\mathrm{CH_4}\), and 0.77\% for \(\mathrm{NH_3}\). Their 95th
percentiles were 0.85, 2.52, and 1.92\%, respectively.

\begin{figure}[htbp]
\centering
\includegraphics[width=\textwidth]{Fig3_campaign1_sampling_configuration_deviation.png}
\caption{Deviation of six reduced vertical sampling configurations from the
full 51-location mean in 743 complete two-hour blocks. Boxes show the
interquartile range and centre lines show medians; whiskers extend to
1.5 times the interquartile range. Green and yellow bands indicate deviations
within \(\pm5\%\) and \(\pm10\%\), respectively. Configurations labelled top,
middle, or bottom contain 17 locations; paired-height configurations contain
34 locations. Outliers are omitted visually but retained in all summaries.}
\label{fig:configurations}
\end{figure}

\begin{table}[htbp]
\centering
\caption{Performance of selected reduced configurations relative to the full
51-location mean. Values are calculated across 743 complete two-hour blocks.
MAD denotes median absolute deviation and \(P_{95}\) denotes the 95th
percentile of absolute deviation.}
\label{tab:configurations}
\begin{tabular}{llrrr}
\toprule
Gas & Configuration & Mean RE (\%) & MAD (\%) & \(P_{95}\) (\%)\\
\midrule
\(\mathrm{CO_2}\) & Top only & 0.68 & 0.82 & 3.71\\
 & Middle only & -0.18 & 0.52 & 1.70\\
 & Bottom only & -0.50 & 0.78 & 3.28\\
 & Top + bottom & 0.09 & 0.26 & 0.85\\
\(\mathrm{CH_4}\) & Top only & 1.01 & 2.42 & 9.75\\
 & Middle only & -0.23 & 1.46 & 5.03\\
 & Bottom only & -0.78 & 2.17 & 8.37\\
 & Top + bottom & 0.12 & 0.73 & 2.52\\
\(\mathrm{NH_3}\) & Top only & 1.18 & 2.34 & 9.67\\
 & Middle only & 0.22 & 1.28 & 3.84\\
 & Bottom only & -1.40 & 2.02 & 8.04\\
 & Top + bottom & -0.11 & 0.64 & 1.92\\
\bottomrule
\end{tabular}
\end{table}

\subsection{Entropy and temporal-resolution sensitivity}

The most temporally informative location depended on the response
(Fig.~\ref{fig:entropy}). The highest normalised entropy occurred at locations
4, 7, and 15 for \(\mathrm{CO_2}\), \(\mathrm{CH_4}\), and
\(\mathrm{NH_3}\), respectively; no single location was universally
informative. Increasing the averaging interval from 2 to 24 h reduced the
median spatial CV from 7.73 to 4.60\% for \(\mathrm{CO_2}\), from 19.11 to
9.56\% for \(\mathrm{CH_4}\), and from 22.88 to 20.11\% for
\(\mathrm{NH_3}\). The persistent ammonia CV indicated that its
heterogeneity was not primarily a short-timescale artefact.

\begin{figure}[htbp]
\centering
\includegraphics[width=\textwidth]{Fig6_june_august_entropy_v03.png}
\caption{Normalised four-bin Shannon entropy of corrected gas concentrations
and gas ratios at each Campaign 1 location. Entropy was calculated across 743
complete two-hour blocks and normalised by \(\log_2(4)\). Numbers identify
sampling locations. Higher values indicate broader and more even occupation
of the four empirical bins, not necessarily a greater mean concentration.}
\label{fig:entropy}
\end{figure}

\subsection{Agreement between full and reduced designs}

The top-and-bottom mean closely reproduced the full 51-location estimate
throughout Campaign 1 (Fig.~\ref{fig:agreement}). Lin's CCC was 0.999 for all
three gases. Mean relative biases were 0.09\% for \(\mathrm{CO_2}\), 0.12\%
for \(\mathrm{CH_4}\), and -0.11\% for \(\mathrm{NH_3}\). The corresponding
RMSE values were \SI{2.69}{ppm}, \SI{0.23}{ppm}, and \SI{0.019}{ppm}.
Agreement for the spatial heterogeneity metric was lower than agreement for
the mean but remained strong, with CCC values between 0.953 and 0.968.

\begin{figure}[htbp]
\centering
\includegraphics[width=\textwidth]{Fig5_june_august_agreement_v03.png}
\caption{Agreement between the full 51-location mean and the reduced
34-location top-and-bottom mean in 743 complete two-hour blocks. Solid lines
represent exact 1:1 agreement. CCC denotes Lin's concordance correlation
coefficient, and bias is the mean relative difference from the full-density
estimate.}
\label{fig:agreement}
\end{figure}

\subsection{Longer-term application of the reduced design}

The reduced-height principle was applied from mid-November to the end of
December 2024 in Campaign 2 (Fig.~\ref{fig:campaign2}). Mean concentrations
remained higher at top than at bottom locations: \(\mathrm{CO_2}\) averaged
approximately 640 and \SI{625}{ppm}, \(\mathrm{CH_4}\) averaged 18.6 and
\SI{17.4}{ppm}, and \(\mathrm{NH_3}\) averaged 0.96 and \SI{0.83}{ppm} at
the top and bottom, respectively. The mean \(\mathrm{CH_4}/\mathrm{CO_2}\)
ratios were 2.75\% and 2.64\%, and the mean
\(\mathrm{NH_3}/\mathrm{CO_2}\) ratios were 0.149\% and 0.131\% at the top
and bottom.

\begin{figure}[htbp]
\centering
\includegraphics[width=\textwidth]{Fig5_campaign2_concentrations_reduced_design.png}
\caption{Campaign 2 concentrations measured with the reduced top-and-bottom
network. Points denote means of daily location means and vertical bars denote
95\% confidence intervals calculated across daily blocks. Orange and
light-blue symbols denote top and bottom locations, respectively.}
\label{fig:campaign2}
\end{figure}

\section{Discussion}

\subsection{Concentration heterogeneity and vertical structure}

The full-scale measurements confirmed that the assumption of homogeneous gas
concentrations was not supported throughout the barn volume. Concentrations
were generally greater at the top than at the bottom, but the gradient varied
among horizontal positions and gases. This gas-specific behaviour agrees with
the vertical variability reported by \citet{Mendes2015} and
\citet{Durso2024}. Those studies showed that vertical profiles cannot be
generalised from one gas to another and that correlations among heights may
remain incomplete. The present network extended this evidence from a limited
number of vertical locations to 17 horizontally distributed triplets.

The differences among locations also support the horizontal-distribution
findings of \citet{Durso2022}, who showed that an individual location was
unlikely to provide a representative barn mean and that both the number and
separation of locations affected variability. In the present barn, the
contrast among the 17 horizontal positions was sufficiently large that
vertical reduction was only successful when spatial coverage across the barn
was retained.

\subsection{Why the top-and-bottom configuration was retained}

The middle-only configuration gave a comparatively small mean bias, but it did
not preserve the vertical extremes. Selection of a network solely because it
reproduces the overall mean would therefore conceal part of the spatial
heterogeneity that the measurement system was intended to characterise. In
contrast, the combined top-and-bottom configuration balanced the positive
deviation at the top against the negative deviation at the bottom while
retaining sensitivity to vertical gradients.

This result is consistent with the general finding that sensor placement is
not interchangeable with sensor number. \citet{DeVogeleer2017} demonstrated
that configurations containing the same number of sensors could yield
different accuracy because of their spatial pattern. \citet{Janke2020}
reported particularly large errors when concentration and velocity were
represented by one height and improved performance from vertically composite
sampling. The present field measurements provide an analogous result for
internal barn concentrations: retaining two vertically contrasting levels was
more representative than retaining either level alone.

The near-unity CCC for the 34-location mean demonstrates that the reduction
introduced little loss of information about the overall concentration.
Nevertheless, agreement for the spatial heterogeneity metric was weaker than
agreement for the mean. This distinction is important because a reduced
network may reproduce the mean through compensating positive and negative
deviations while slightly changing its description of spatial spread.
Accordingly, both outcomes should be evaluated when sampling networks are
optimised.

\subsection{Implications for sensor-network optimisation}

The reduction from 51 to 34 locations lowered the number of internal sampling
points by one third without materially changing the two-hour barn mean.
Campaign 2 subsequently showed that the reduced-height principle remained
operational over a longer measurement period. This progression follows the
sensor-network concept advanced by \citet{Chen2026}: a full network can first
be used to identify informative spatial structure, after which a smaller
network can be retained for routine monitoring.

The design should not, however, be transferred directly to barns with
different geometry or ventilation without validation. Optimal positions are
expected to depend on openings, internal obstructions, source distribution,
fans, wind direction, and the intended response variable. The present
findings support a general procedure---dense initial mapping followed by
objective reduction---rather than a universal requirement for exactly 34
locations.

\subsection{Limitations and future research}

Several limitations should be considered. First, Campaign 1 was conducted in
summer and Campaign 2 in late autumn and winter; consequently, the complete
range of seasonal ventilation conditions was not represented. Second, the
four analysers were assigned to different location ranges. Although the
preliminary co-located comparison was used to harmonise FTIR and CRDS values,
analyser type and barn location were not completely separable during Campaign
1. Third, sequential switching meant that all locations were not sampled at
the exact same second. Two-hour complete blocks were therefore used to compare
configurations. Fourth, the 51-location estimate was treated as the
full-density reference rather than as an absolute true barn concentration.

The Campaign 2 network contained 32 rather than 34 locations because two
locations were removed for identified site-specific problems. Therefore,
Campaign 2 provided a longer-term application of the reduced principle but
not an independent 51-versus-34 validation. Future investigations should
combine simultaneous meteorological, animal, and barn-climate measurements
with independent validation of strategically reduced networks.

\section{Conclusions}

Gas concentrations in the investigated naturally ventilated dairy barn were
heterogeneous in both the vertical and horizontal directions. Measurements
from only the top systematically overrepresented the full-density mean,
whereas measurements from only the bottom systematically underrepresented it.
The middle height reproduced the overall mean more closely but did not retain
the observed vertical extremes.

A spatially distributed top-and-bottom network preserved the mean
\(\mathrm{CO_2}\), \(\mathrm{CH_4}\), and \(\mathrm{NH_3}\)
concentrations obtained from the full 51-location design while reducing the
number of internal locations by one third. Dense initial mapping followed by
quantitative configuration testing provides a defensible route for reducing
sampling effort without assuming homogeneous mixing.

\section*{Acknowledgements}

The technical and animal-care staff of the experimental dairy facility are
gratefully acknowledged. \textit{[Funding project names, grant numbers, and
specific technical contributors are to be confirmed before submission.]}

\section*{Author contributions}

Harsh Sahu: conceptualisation, methodology, investigation, data curation,
formal analysis, visualisation, and writing--original draft. Long Chen:
methodology, formal analysis, and writing--review and editing. Thomas Amon:
conceptualisation, supervision, funding acquisition, and writing--review and
editing. Tariq Mehmood: investigation and writing--review and editing. Aditya
Rawat: investigation and writing--review and editing. David Janke:
conceptualisation, methodology, supervision, project administration, and
writing--review and editing.

\section*{Declaration of competing interest}

The authors declare that they have no known competing financial interests or
personal relationships that could have appeared to influence the work
reported in this paper.

\section*{Data availability}

The processed data and analysis code will be made available in a public
repository upon acceptance. \textit{[Repository and DOI to be added.]}

\bibliographystyle{apalike}
\bibliography{references}

\end{document}
